% ED26-SCI-GoldbergOracle-001. Source: Sun Studio 11, 819-3693-10, Appendix D.
\chapter{Numerical Error and Computed Conclusions}
\label{sec:p12-numerical-error-and-inference}

A computed number supports a sign, an ordering or an accuracy claim only
when the accompanying error bound supports that claim. This chapter follows
the rounding-error presentation in the November 2005 \emph{Numerical
Computation Guide}, Appendix D: an edited reprint of Goldberg's article,
with its source-specific arithmetic assumptions retained
\cite[Sections D.3, D.6.1 and D.10.1--D.10.2]{GoldbergOracle8193693}.
It does not assert that a particular processor, compiler or library satisfies
those assumptions. The SMT interface in \cref{def:p12-smt-problem} concerns
satisfaction in a declared theory. Numerical stability additionally compares
a computed result with a specified real-valued target.

\section{Values, representation and error}

\begin{hbdefinition}{def:p12-numerical-format}{A finite radix format and its idealized interior}
Fix an even integer base \(\beta\geq2\), an integer precision \(p\geq2\),
and integer exponent limits \(e_{\min}\leq e_{\max}\). The nonzero normal values are
\[
 \pm\left(\sum_{j=0}^{p-1}d_j\beta^{-j}\right)\beta^e,
 \quad 1\leq d_0<\beta,\quad 0\leq d_j<\beta,
 \quad e_{\min}\leq e\leq e_{\max},
\]
where the digits and exponent are integers. Include zero separately.
The coefficient is the \emph{significand}; multiplication by \(\beta^e\)
in this notation is an exact real scaling, not a computed multiplication.
The proofs below use the same significands with unrestricted integer
exponents. For transfer to a finite implementation, assume its executed
operations follow the declared rounding procedure, that their nonzero
stored inputs and outputs are normal, and that the magnitude of each
nonzero exact real operation result stays between its smallest and largest positive normal
magnitudes. The exact real result need not itself be a format value:
rounding is precisely what handles that discrepancy. Rescalings used
only inside a proof are not executed operations. These sufficient range
conditions exclude overflow, underflow, exceptional values and
flush-to-zero behavior from the relative-error analysis.
\index{floating-point format}\index{significand}\index{rounding error}
\end{hbdefinition}

\begin{hbdefinition}{def:p12-numerical-rounding}{Nearest rounding, ulps and relative error}
Write \(R(z)\) for rounding a real value to a nearest value of the
unrestricted-exponent format, with \(R(0)=0\). At a midpoint choose the
candidate whose least significant retained digit is even. This is a
deterministic tie rule, not a statement about the frequency of ties.
For a nonzero real \(z\) with \(\beta^e\leq |z|<\beta^{e+1}\), set
\[
 \operatorname{ulp}(z)=\beta^{e-p+1},\qquad
 u=\tfrac12\beta^{1-p}.
\]
For an approximation \(\widehat z\), the absolute error is
\(|\widehat z-z|\), and its relative error, when \(z\ne0\), is
\(|\widehat z-z|/|z|\). The interval
\([\beta^e,\beta^{e+1})\) is a \emph{radix interval}, called a
binade when \(\beta=2\). An error in ulps must state which value's radix interval
sets the unit. Goldberg calls \(u\) machine epsilon; the distance from
\(1\) to the next larger normal value is \(2u\), a different quantity.
At a zero target relative error is undefined, so use an absolute bound.
\index{unit in the last place}\index{relative error}\index{round to even}
\end{hbdefinition}

\begin{hbtheorem}{thm:p12-numerical-nearest}{proved=annex:p12.numerical.nearest}{Spacing and the nearest-rounding bound}
For \(z\ne0\),
\[
 |R(z)-z|\leq\tfrac12\operatorname{ulp}(z),\qquad
 R(z)=z(1+\delta),\quad |\delta|\leq u.
\]
In a radix interval, a fixed half-ulp absolute error corresponds to relative
errors between \(u/\beta\) and \(u\); this factor-\(\beta\) variation
is the source's \emph{wobble}. A finite normal-only format has no such
uniform relative guarantee arbitrarily close to zero.
\end{hbtheorem}

\begin{hbexample}{def:p12-numerical-format-examples}{Finite precision does not form a real field}
For \(\beta=2,p=3,e_{\min}=-1,e_{\max}=2\), the positive normal
values have significands \(1,1.25,1.5,1.75\), each multiplied by one of
\(2^{-1},1,2,4\). These are sixteen values; their negatives and zero
give thirty-three real values. The real number \(1/10\) has no finite
binary expansion: a finite binary fraction has a power-of-two denominator
after reduction, whereas \(1/10\) has a factor of five.
In decimal precision three, \(R(3.14159)=3.14\); the absolute error
\(0.00159\) is \(0.159\) ulp in that radix interval.
With sufficiently wide decimal exponent limits,
\[
 R(R(1000+1)-1000)=0,\qquad R(1000+R(1-1000))=1.
\]
The first sum rounds to \(1000\); the second inner difference is
\(-999\). Parentheses therefore affect the computed answer.
\end{hbexample}

\section{Guard digits and exact subtraction}

\begin{hbdefinition}{def:p12-numerical-guard}{The arithmetic procedures being compared}
To add or subtract two nonnegative format values, align their radix
points at the exponent of the larger operand. Retain the leading \(p+g\)
positions, counting the leading position before the radix point; discard
the smaller operand's digits beyond them. Perform that aligned operation
exactly, then round its result to \(p\) significant digits. Here \(g=0\)
means no guard digit and \(g=1\) means one guard digit. Signed operations
reduce to a same-sign sum or a difference of magnitudes.
Write \(x\oplus y\) and \(x\ominus y\) for the declared computed
sum and difference. In this chapter these use at least one guard digit
unless a no-guard example is explicitly named. Write
\(x\otimes y=R(xy)\) and \(x\oslash y=R(x/y)\), with \(y\ne0\)
for division. An \emph{exactly rounded} addition or subtraction instead
means \(R(x+y)\) or \(R(x-y)\) directly. One guard digit does not
in general implement that stronger procedure.
\end{hbdefinition}

\begin{hbtheorem}{thm:p12-numerical-no-guard}{proved=annex:p12.numerical.no-guard}{The source's no-guard failure}
In every such base and precision there are positive format values for
which subtraction with no guard digit has relative error \(\beta-1\).
In decimal precision three, the no-guard computation of \(10.1-9.93\)
returns \(0.2\), while its real value is \(0.17\).
\end{hbtheorem}

\begin{hbtheorem}{thm:p12-numerical-guard-bound}{proved=annex:p12.numerical.guard-bound}{One guard digit and same-sign addition}
For positive format values \(x\ne y\), one-guard subtraction satisfies
\[
 \frac{|(x\ominus y)-(x-y)|}{|x-y|}
 <(1+2/\beta)u\leq2u.
\]
For nonnegative \(x,y\) with \(x+y>0\), same-sign addition has relative
error at most \(2u\), even with no guard digit. Consequently the signed
one-guard addition/subtraction model has relative error at most \(2u\)
at every nonzero exact result. A zero exact result is computed as zero.
These are the source's Theorems 2, 9 and 10 with the zero case explicit.
\end{hbtheorem}

\begin{hbtheorem}{thm:p12-numerical-close-subtraction}{proved=annex:p12.numerical.close-subtraction}{Exact difference of nearby format values}
If \(x,y>0\) are format values and \(y/2\leq x\leq2y\), then
\(x-y\) is a format value and one-guard or exactly rounded subtraction
returns it exactly. The finite-range version requires representability
of this difference; the factor-two condition alone does not authorize
discarding a subnormal result.
\end{hbtheorem}

\begin{hbexample}{def:p12-numerical-guard-not-nearest}{A guard digit can differ from nearest rounding}
In decimal precision three, aligning \(110-8.59\) with one guard digit
gives \((1.100-0.085)10^2=101.5\), which rounds to \(102\).
The exact difference is \(101.41\), whose nearest value is \(101\).
The guard result's relative error is \(0.59/101.41>u=0.005\),
while remaining below \(2u\).
\end{hbexample}

\section{Cancellation and algebraically equivalent expressions}

\begin{hbtheorem}{thm:p12-numerical-product-errors}{proved=annex:p12.numerical.product-errors}{Combining multiplicative errors}
If \(|\delta_i|\leq a_i<1\), then
\[
 \left|\prod_{i=1}^n(1+\delta_i)-1\right|
 \leq\prod_{i=1}^n(1+a_i)-1.
\]
In particular, for exactly specified format inputs \(x,y\) with
\(x^2\ne y^2\), evaluation as
\((x\ominus y)\otimes(x\oplus y)\) has relative error at most
\[
 (1+2u)^2(1+u)-1=5u+8u^2+4u^3.
\]
This bound retains all terms. It concerns error relative to \(x^2-y^2\),
not to an unspecified pair of real quantities approximated by \(x,y\).
\end{hbtheorem}

\begin{hbtheorem}{thm:p12-numerical-cancellation}{proved=annex:p12.numerical.cancellation}{Exact subtraction cannot repair inaccurate operands}
Suppose \(\widehat A=A+e_A\) and \(\widehat B=B+e_B\), with
\(|e_A|\leq E_A\), \(|e_B|\leq E_B\). If their subtraction is exact,
its absolute error relative to \(A-B\) is at most \(E_A+E_B\).
When \(A\ne B\), the corresponding relative bound is
\((E_A+E_B)/|A-B|\), which need not be small. If subtraction itself
has absolute error at most \(E_S\), add \(E_S\) to that budget.
\end{hbtheorem}

\begin{hbexample}{def:p12-numerical-discriminant}{An exact difference with contaminated input products}
Use decimal precision three and exact inputs
\(a=1.22,b=3.34,c=2.28\). Their real discriminant is
\[
 b^2-4ac=11.1556-11.1264=0.0292.
\]
The nearest-rounded products are \(11.2\) and \(11.1\).
Subtracting them gives \(0.1\) exactly, but its absolute error is
\(0.0708\), or \(70.8\) ulps measured at \(0.100\).
Its relative error is \(177/73\). The subtraction is accurate for
the stored operands; the computed discriminant is inaccurate for the
stated real expression.
\end{hbexample}

\begin{hbtheorem}{thm:p12-numerical-quadratic}{proved=annex:p12.numerical.quadratic}{Conjugate evaluation of a quadratic root}
Let \(a\ne0\), \(b>0\), and \(d=b^2-4ac\geq0\). Then
\[
 \frac{-b+\sqrt d}{2a}=-\frac{2c}{b+\sqrt d}.
\]
For \(b<0\), the corresponding conjugate form avoids cancellation in
the other numerator. The identity removes that numerator's cancellation;
it neither proves that the computed discriminant is accurate nor bounds
uncertainty in the input coefficients.
\end{hbtheorem}

\begin{hbexample}{def:p12-numerical-quadratic-trace}{A root lost by direct evaluation}
For \(a=c=1,b=100\) in decimal precision three, exactly rounded
evaluation gives \(R(10000-4)=10000\), then \(R(\sqrt{10000})=100\).
The direct small-root numerator becomes zero. The conjugate evaluation
returns \(-2/200=-0.0100\). The actual root is
\(-2/(100+\sqrt{9996})\), approximately \(-0.0100010002\).
The conjugate computation retains the small root in this example;
its reliability follows from the arithmetic and bounds, not merely from
printing additional decimal places.
\end{hbexample}

\section{A stable triangle-area evaluation}

\begin{hbdefinition}{def:p12-numerical-triangle}{The ordered-side expression}
Let \(a\geq b\geq c>0\) and \(a\leq b+c\). Put
\[
 q=(a+(b+c))(c-(a-b))(c+(a-b))(a+(b-c)).
\]
Its computed version forms exactly these four parenthesized factors,
then multiplies them with three nearest-rounded multiplications to
obtain \(Q\). The computed area is
\(\operatorname{SQRT}(Q)\oslash4\). Assume
\(\operatorname{SQRT}(0)=0\), and for \(t>0\) assume that
\(\operatorname{SQRT}(t)\) is a format value within
\(\tfrac12\operatorname{ulp}(\sqrt t)\) of \(\sqrt t\).
No particular tie direction for this library result is needed.
Keep the parentheses; algebraic reassociation changes the computation.
\end{hbdefinition}

\begin{hbtheorem}{thm:p12-numerical-triangle-radicand}{proved=annex:p12.numerical.triangle-radicand}{The complete radicand error bound}
If the sides are format values, the triangle is nondegenerate, and
\(0<u\leq0.005\), then \(|Q/q-1|<16u\).
For a degenerate triangle \(a=b+c\), the same procedure returns
\(Q=0\) exactly under the stated range assumptions.
\end{hbtheorem}

\begin{hbtheorem}{thm:p12-numerical-triangle-area}{proved=annex:p12.numerical.triangle-area}{The source's eleven-unit area guarantee}
The real area is \(\sqrt q/4\). In the nondegenerate case,
\[
 \left|\frac{\operatorname{SQRT}(Q)\oslash4}{\sqrt q/4}-1\right|
 <11u
\]
under the preceding assumptions. In binary, division by four is exact
when the scaled value stays in range; the displayed bound also allows
the extra nearest rounding required in another even base.
The degenerate case returns area zero exactly.
\end{hbtheorem}

\begin{hbexample}{def:p12-numerical-flat-triangle}{A complete flat-triangle trace}
For \(a=9.00,b=c=4.53\), the real semiperimeter is \(9.03\),
and the area is \(\sqrt{5.485725}\), approximately \(2.34216\).
In decimal precision three, the order
\(R(R(b+c)+a)/2\) gives \(R(9.06+9.00)/2=18.1/2=9.05\).
Heron's four factors are then \(9.05,0.05,4.52,4.52\).
Multiplying left to right with nearest-even rounding gives
\(0.452,2.04,9.22\), and the rounded square root is \(3.04\).
For the ordered-side expression, the four computed factors are
\(18.1,0.0600,9.00,9.00\). The three products are
\(1.09,9.81,88.3\); the rounded square root is \(9.40\), and
division by four returns \(2.35\). The evaluation order is part
of each calculation, not an unmentioned implementation detail.
\end{hbexample}

\section{Recovering a small logarithmic increment}

\begin{hbdefinition}{def:p12-numerical-log-increment}{The compensated logarithm expression}
For a format value \(0\leq x<3/4\), compute \(z=1\oplus x\).
If \(z=1\), return \(x\). Otherwise compute
\[
 t=z\ominus1,\qquad
 L=x\otimes\bigl(\operatorname{LN}(z)\oslash t\bigr),
\]
Here \(\operatorname{LN}(z)\), when \(z>1\), is assumed to be a
format value within \(\tfrac12\operatorname{ulp}(\ln z)\) of
\(\ln z\). No particular tie direction for this library result is
needed. No rounding guarantee for logarithms follows
just from an assumption about addition or multiplication.
\end{hbdefinition}

\begin{hblemma}{thm:p12-numerical-mu}{proved=annex:p12.numerical.mu}{The slowly varying logarithmic ratio}
Define \(\mu(0)=1\) and \(\mu(x)=\ln(1+x)/x\) for \(x>0\).
For \(x\geq0\),
\[
 \mu(x)=\int_0^1\frac{dt}{1+xt},\qquad
 0<\mu(x)\leq1,\qquad
 -\tfrac12\leq(\ln\mu)'(x)\leq0.
\]
Also \(|\mu'(x)|\leq1/2\), and \(\mu(x)\geq1/2\) on
\([0,1]\). These statements supply the bounds needed even if rounding
moves \(x\) slightly above \(3/4\).
\end{hblemma}

\begin{hbtheorem}{thm:p12-numerical-log-increment}{proved=annex:p12.numerical.log-increment}{The source's five-unit logarithm guarantee}
If \(0<x<3/4\), \(u<0.1\), addition and subtraction use at least one
guard digit, multiplication and division are nearest-rounded, and the
stated half-ulp logarithm assumption holds, the procedure's relative
error with respect to \(\ln(1+x)\)
is at most \(5u\). For \(x=0\) the returned value is zero exactly.
\end{hbtheorem}

\begin{hbexample}{def:p12-numerical-small-increment}{A lost increment and its downstream use}
In binary precision ten, let \(x=2^{-12}\). Then
\(1\oplus x=1\), so evaluating a logarithm after that addition
returns \(\ln1=0\). The compensated branch returns \(x\), with
the nonzero-target relative guarantee just proved.
For \(x>0\) and an integer \(n\geq1\), the exact geometric sum is
\[
 \sum_{j=0}^{n-1}(1+x)^j=\frac{(1+x)^n-1}{x}
 \quad(x\ne0),
\]
and \((1+x)^n=\exp(n\ln(1+x))\). Thus an error in the stored
increment can be amplified by exponentiation. The source's example
uses \(n=365,x=0.06/365\) and multiplies this sum by one hundred;
its real value is approximately \(37614.0473\)
\cite[Section D.3.4, pp.~D-11--D-12]{GoldbergOracle8193693}.
These identities justify using an accurate logarithmic increment;
they do not remove the separate rounding of the exponential,
subtraction or division in a complete implementation.
\end{hbexample}

\section{Exact rounding and retained information}

\begin{hbtheorem}{thm:p12-numerical-no-drift}{proved=annex:p12.numerical.no-drift}{The source's round-to-even no-drift result}
For format values \(x,y\), let \(x_0=x\) and
\[
 x_{j+1}=R\bigl(R(x_j-y)+y\bigr).
\]
In the source's exactly rounded, round-to-even model, with no range
loss, either every \(x_j=x\), or every \(x_j=x_1\) for \(j\geq1\).
This is the source's Theorem 5, attributed there to Reiser and Knuth.
\end{hbtheorem}

\begin{hbexample}{def:p12-numerical-drift-example}{Changing only the midpoint rule}
In decimal precision three take \(x=1.00,y=-0.555\).
Rounding positive midpoints upward gives
\(1.00+0.555\mapsto1.56\), then
\(1.56-0.555\mapsto1.01\). Repeating gives \(1.02\) next.
Round to even instead gives \(1.56\), followed by \(1.00\),
and the same pair repeats. The example concerns a deterministic
midpoint convention; it provides no model of rounding errors as
independent random variables.
\end{hbexample}

\begin{hbdefinition}{def:p12-numerical-coarse-rounding}{Rounding away a block of low digits}
For a positive format value, scale exactly to an integer
\(X\in[\beta^{p-1},\beta^p)\). To discard \(k\) low digits,
where \(0<k<p\), set \(B=\beta^k\), round \(X/B\) to a nearest
integer with ties to the even integer, multiply by \(B\), and undo
the scaling. Denote this value by \(H_k(x)\), extend by sign symmetry,
and set \(H_k(0)=0\). The definition also covers a carry into a new
leading digit; it specifies the tie convention even when only one
significant digit remains. Write \(\operatorname{nint}\) for the
nearest-integer operation used here.
\end{hbdefinition}

\begin{hbtheorem}{thm:p12-numerical-split-rounding}{proved=annex:p12.numerical.split-rounding}{The exact split-rounding identity, including a carry}
With all three operations exactly rounded, set
\[
 m=\beta^k+1,\quad T=R(mx),\quad S=R(T-x),\quad h=R(T-S).
\]
Then \(h=H_k(x)\). The statement assumes no range loss at any step;
in particular \(mx\) must not overflow a finite implementation.
\end{hbtheorem}

\begin{hbtheorem}{thm:p12-numerical-split}{proved=annex:p12.numerical.split}{Two short components and exact partial products}
Let \(k=\lceil p/2\rceil\), and require \(p\) even when
\(\beta>2\). Compute \(h\) by the preceding identity and put
\(l=R(x-h)\). Then
\[
 x=h+l
\]
exactly, and each nonzero component has at most \(\lfloor p/2\rfloor\)
significant radix digits. If \(y=h_y+l_y\) is split in the same way,
each of the four products in
\[
 xy=h h_y+h l_y+l h_y+l l_y
\]
is exactly representable at precision \(p\), under the range assumptions.
Their real sum is the product; an arbitrary rounded summation order
need not preserve that sum exactly.
\end{hbtheorem}

\begin{hbexample}{def:p12-numerical-split-example}{A product difference recovered from short components}
In decimal precision four, \(b=3.476,a=3.463,c=3.479\) split as
\[
 b=3.5-0.024,\quad a=3.5-0.037,\quad c=3.5-0.021.
\]
The exact expansions are
\(b^2=12.25-0.168+0.000576\) and
\(ac=12.25-0.2030+0.000777\). Subtracting corresponding terms and
then rounding gives \(0.0350-0.000201\), rounded to \(0.03480\).
Directly rounding the products gives \(12.08-12.05=0.03\).
The exact difference is \(0.034799\).
In binary precision three, \(x=7,k=2,m=5\) instead illustrates the
rounding assumption: \(R(35)=32\). One-guard subtraction of seven
from thirty-two returns twenty-eight, giving components four and three;
three is not a one-bit component. Exactly rounded subtraction returns
twenty-four, giving the correct components eight and minus one.
\end{hbexample}

\begin{hbtheorem}{thm:p12-numerical-round-trip}{proved=annex:p12.numerical.round-trip}{An even-radix divide-and-multiply identity}
Let \(m,n\) be exactly representable integers in the declared even base,
\(|m|<\beta^{p-1}\), and \(n=\beta^i+\beta^j>0\) for nonnegative integers
\(i,j\). If division and multiplication are exactly rounded and
all needed values remain in range, then
\[
 R\bigl(n R(m/n)\bigr)=m.
\]
The binary case is the source's Theorem 7; the proof also establishes
the even-radix extension mentioned immediately after it.
This is an exact identity for the declared operands; it is not an error
bound for unknown real quantities those integers might encode.
\end{hbtheorem}

\begin{hbtheorem}{thm:p12-numerical-root-sensitivity}{proved=annex:p12.numerical.root-sensitivity}{A precise discriminant-to-root error bound}
For \(d,\widehat d\geq0\) with \(|d-\widehat d|\leq E\),
\[
 |\sqrt{\widehat d}-\sqrt d|\leq\sqrt E.
\]
If both are at least \(d_0>0\), the sharper bound
\(E/(2\sqrt{d_0})\) holds. With \(a\ne0,b\) fixed, either exact
quadratic-root expression consequently changes by at most
\(\sqrt E/(2|a|)\) when only its discriminant changes.
Additional rounding of the square root or root arithmetic needs its
own error terms. A negative computed discriminant is not covered by
these real-square-root bounds.
\end{hbtheorem}

\section{From error bounds to a warranted conclusion}

\begin{hbtheorem}{thm:p12-numerical-certification}{proved=annex:p12.numerical.certification}{Sign, ordering and accuracy certificates}
Suppose a real target \(v\) and computed value \(\widehat v\) satisfy
\(|v-\widehat v|\leq E\), where \(E\geq0\) is a proved bound.
Then \(\widehat v-E>0\) certifies \(v>0\), and
\(\widehat v+E<0\) certifies \(v<0\). If the interval contains zero,
these bounds alone certify neither strict sign. Two targets satisfy
\(v<w\) whenever
\(\widehat v+E_v<\widehat w-E_w\).
If \(|\widehat v|>E\), relative error is bounded by
\[
 \frac{|\widehat v-v|}{|v|}\leq
 \frac{E}{|\widehat v|-E}.
\]
An interval whose endpoint is exactly zero does not certify a strict sign.
\end{hbtheorem}

\begin{hbtheorem}{thm:p12-numerical-data-budget}{proved=annex:p12.numerical.data-budget}{Separate input uncertainty from arithmetic error}
For exact stored coefficients \(a,b,c\), suppose the target coefficients
obey \(|a_*-a|\leq\eta_a\), \(|b_*-b|\leq\eta_b\),
\(|c_*-c|\leq\eta_c\). Then
\[
 |(b_*^2-4a_*c_*)-(b^2-4ac)|
 \leq 2|b|\eta_b+\eta_b^2
 +4\bigl(|a|\eta_c+|c|\eta_a+\eta_a\eta_c\bigr).
\]
Add this quantity to any proved arithmetic-error bound for the stored
coefficient calculation. More accurate arithmetic reduces the latter
term; it does not reduce the stated uncertainty in the inputs.
\end{hbtheorem}

\begin{hbexample}{def:p12-numerical-certification-cases}{A certified sign and two limits of certification}
In the conjugate-root example let \(s=\sqrt{9996}\), so
\(99<s<100\). The exact error of the returned value \(-0.0100\) is
\[
 \frac{2}{100+s}-\frac1{100}
 =\frac{4}{100(100+s)^2}
 <E=\frac{4}{3960100}.
\]
The resulting interval lies strictly below zero, certifying a negative
root. The general relative-error certificate is
\(E/(0.01-E)=4/39597\), approximately \(0.000101018\).

In the earlier discriminant trace, nearest rounding alone gives
\(|A-\widehat A|\leq u|\widehat A|/(1-u)\) for each product:
rearrange \(|\widehat A-A|\leq u|A|\) and use
\(|A|\leq|\widehat A|/(1-u)\). Thus the exactly subtracted
\(\widehat d=0.1\) has the generic absolute budget
\[
 E_d=\frac{0.005}{0.995}(11.2+11.1)=\frac{223}{1990}>0.1.
\]
That interval contains zero and cannot certify the sign, despite the
positive displayed number. Exact input arithmetic happens to establish
the positive value \(0.0292\); it is stronger information than this
generic budget.

For an ordering example, certificates \(v\in[0.008,0.012]\) and
\(w\in[0.009,0.013]\) do not determine \(v<w\), even if the
reported centers are \(0.010\) and \(0.011\). Both
\((v,w)=(0.008,0.013)\) and \((0.012,0.009)\) satisfy them.
Input uncertainty enlarging a numerical enclosure must be included
before any of these decisions.
\end{hbexample}

\begin{hbmethod}{def:p12-numerical-inference-procedure}{Using a numerical value as evidence}
State the real target and the stored inputs separately. Declare the format,
rounding of each operation, evaluation order and exceptional-case policy.
Check the range assumptions. Obtain an arithmetic-error bound for that
specific procedure; propagate input uncertainty separately. Apply the
sign, ordering or relative-accuracy certificate to the combined interval.
When intervals overlap, obtain sharper bounds or better input information,
or retain the unresolved conclusion. A solver certificate for a floating-point
formula establishes a result in its specified arithmetic semantics; the
real-target comparison still requires this error argument.
\index{numerical stability}\index{sign certification}\index{input uncertainty}
\end{hbmethod}
